% PDF 排版微调（由 build_pdf.sh 通过 pandoc -H 引入）
\usepackage{etoolbox}
\usepackage{graphicx}

% 图片统一缩放到版心宽度，保持纵横比
\setkeys{Gin}{width=\linewidth,keepaspectratio}

% 表格自动缩小字号（7.8pt：既压住宽度，又高于位图字体的 7pt 门槛）
% 列表前保证至少 5 行空间，避免"图 X 结论"标签孤立在页底
\usepackage{needspace}
\AtBeginEnvironment{itemize}{\needspace{5\baselineskip}}
\AtBeginEnvironment{longtable}{\fontsize{7.8}{9.2}\selectfont}
\setlength{\tabcolsep}{3.5pt}
\renewcommand{\arraystretch}{1.05}
\setlength{\LTpre}{6pt}
\setlength{\LTpost}{6pt}

% 代码块字号与换行
\AtBeginEnvironment{Shaded}{\small}

% 标题层级间距收紧
\usepackage{titlesec}
\titlespacing*{\section}{0pt}{10pt}{4pt}
\titlespacing*{\subsection}{0pt}{8pt}{3pt}
\titlespacing*{\subsubsection}{0pt}{6pt}{2pt}


% 中文段落微调
\setlength{\parindent}{0pt}
\setlength{\parskip}{4pt plus 1pt}

% 允许页面内容轻微宽松，减少 overfull 警告
\sloppy

% 图注
\usepackage{caption}
\captionsetup{font=small,labelfont=bf,labelformat=empty,skip=3pt}

% 宽幅多面板图可旋转为横向整页
\usepackage{rotating}
